\section{未来路线：Linear 与 Sparse 会不会合流}

本章把前面几条路线合起来看。节目中提到一个未来理想方案：把 Linear Attention 和 Sparse Attention 结合。直觉上，Linear/Hybrid 负责降低 KV cache 和大部分层的递推成本，Sparse 负责在需要全局检索时选中少数关键 token。理论上，如果 sparse 选得准，它可以替代一部分 Full Attention 层。

\subsection{为什么组合路线有吸引力}

本节解释为什么未来很可能不是单一路线胜出。Linear、Sparse、Full 各自解决不同瓶颈，真正的工程系统往往会把它们组合起来。

Linear Attention 擅长省状态和降低大多数 token 的成本，但可能牺牲一些复杂全局交互；Sparse Attention 保留“看历史中重要位置”的能力，但需要一个足够可靠的 indexer。组合路线希望两者互补：大部分时间用 cheap recurrence，必要时用 sparse retrieval 补全长程信息。

\begin{knowledgebox}{组合路线的三种可能}
\begin{tabularx}{\textwidth}{p{0.24\textwidth}p{0.34\textwidth}X}
\toprule
组合方式 & 机制 & 主要风险 \\
\midrule
Linear + Full & 多数线性层，少数全局层 & 全局层仍可能限制 context window。 \\
Sparse + Full & 多数稀疏层，少数全局层 & KV cache 仍大，indexer 要可靠。 \\
Linear + Sparse & 线性状态加稀疏检索 & 训练稳定性和实现复杂度更高。 \\
\bottomrule
\end{tabularx}
\end{knowledgebox}

\subsection{什么叫“选得准”}

Sparse Attention 的核心不是少看 token，而是少看时还要看对。对于多跳推理、长文档检索、代码依赖和 agent 轨迹，关键信息可能出现在很远的位置。Indexer 如果选错，计算省了，但模型答错；如果选太多，又失去稀疏优势。

\begin{warningbox}{Sparse 的失败模式}
稀疏注意力最危险的失败不是速度慢，而是静默漏掉关键 token。长上下文任务里，漏掉一个早期条件或远处约束，可能导致后续推理全错。
\end{warningbox}

\subsection{本章小结}

未来 Attention 可能不是 Linear、Sparse、Full 三选一，而是组合。真正的难点在于如何让组合路线同时保持表达力、低成本、硬件亲和和训练稳定性。

